\chapterpage{spatial}{比较不同时刻的空间分布}
进入\ui{空间图}（\keys{Ctrl+1}），可以在一张图上比较同一个变量在多个输出时刻的分布。“时刻选择”中，每一行包含从零开始的记录编号，以及以 ns 表示的实际输出时间。
\begin{enumerate}
\item 打开 \code{quickstart.case}，选择 \ui{T}、\ui{XC}、\ui{um} 和 \ui{eV}。单独选中时刻 \ui{\#0}，将样式设为蓝色实线，应用到当前选择。
\item 用同样的方法给 \ui{\#4} 指定橙色虚线，给 \ui{\#8} 指定紫色点划线。按住 Ctrl 选中以上三个时刻，勾选“显示图例”，点击\ui{绘制}。
\end{enumerate}
\Figure{spatial}{在 0、2、4 ns 三个时刻，峰值相同，峰所在的位置不同。根据图例中的时间可以辨认每一条曲线。}{110mm}
长列表可用“全选”或“反选”批量操作。“删除选中”会把相应时刻标成排除状态。需要恢复时，使用“取消删除”或恢复按钮，再选中并重新绘制。选取与排除作用于当前内存中的绘图列表，原始数据文件会保留。
\checknote{运动网格}{绘图时，程序会读取各时间步自身的 \code{XC} 或 \code{X}。空间图 CSV 只有一列位置坐标，各时刻共用该列。第~\pageref{checkcsv}~页说明了这种差异如何影响定量比较。}

\chapterpage{temporal}{读取时间演化曲线}
进入\ui{时间图}（\keys{Ctrl+2}）。全局变量在每个输出时刻有一个存储值。网格模式则从一个空间场量中抽取指定单元，每个选中的单元对应一条曲线。
\begin{enumerate}
\item 选择“全局时间变量”和 \ui{ENIE}，将时间单位设为 \ui{ns}，点击\ui{绘制}。在 0、2、4 ns，示例的读数分别为 10、14、18。
\item 切换到网格单元的时间演化模式，选变量 \ui{T}、单位 \ui{eV} 和 \ui{Cell \#10}，勾选图例并绘制。若要比较三个位置，可按住 Ctrl 再选 \ui{\#6} 和 \ui{\#14}，重新绘制。
\end{enumerate}
\Figure{temporal}{第 \#10 个单元在九个输出时刻的温度。练习网格保持静止，所以这个单元始终位于 50 µm。}{110mm}
多条曲线可按单元指定颜色和线型，先应用样式，再绘图。“按材料着色”需要文件提供材料标识 \code{MID}。练习包没有材料标识。全局模式的线条样式由当前绘图控件决定。
\checknote{读表核对}{本页“导出数据”会写出一个时间列，并为每个已启用且选中的单元写一列。2 ns 对应的 \ui{Cell \#10} 数值为 8 eV。全局模式的第二列用标量名称作为列名。}

\chapterpage{heatmap}{绘制时空热图}
热图中的每一种颜色对应某个时刻、某个位置上的场量值。先沿用空间图和时间图中的变量、坐标与单位，便于用同一数值核对颜色条。
\begin{enumerate}
\item 按 \keys{Ctrl+3} 进入\ui{热图}。选 \ui{T}、\ui{XC}、\ui{um}、\ui{ns} 和 \ui{eV}。在坐标方向中将时间放在横轴，位置放在纵轴。
\item 选择色表 \ui{viridis}，点击颜色范围旁的“自动”，让程序填写下限和上限。保持线性颜色刻度，点击\ui{绘制}。
\end{enumerate}
\Figure{heatmap}{横向为时间，纵向为位置。亮带在 2 ns 左右经过 50 µm，右侧颜色条给出相应的温度，单位为 eV。}{110mm}
在“位置标记”中输入 \code{50}，在“时间标记”中输入 \code{2}，再次绘图即可标出这两个坐标。多个数值用英文逗号分隔。即使交换横纵轴，两个输入框仍分别填写位置和时间。
\checknote{先看轴再点选}{本图的 2 ns、50 µm 对应画布坐标 \((2,50)\)。横轴改为位置、纵轴改为时间后，同一点的画布坐标变为 \((50,2)\)。}

\nextpage{sampling}{颜色范围、缩放与插值}
\subhead{让两张热图可以直接比较}
\begin{enumerate}
\item 用 eV 绘制示例的 \ui{T} 热图。先点击“自动”了解数据范围，然后手动将颜色最小值、最大值设为 2 和 8，再次绘图。
\item 打开待比较的数据，选择相同的变量、坐标、单位、色表和颜色上下限。比较持续时间或位置范围时，也应固定相同的横纵轴范围。
\item 如果物理量均为正且跨越较大数量级，可以使用对数颜色，手动设置大于零且小于上限的下限。需要辨别过零点或正负变化时，使用线性颜色。
\end{enumerate}
“X 范围”和“Y 范围”改变画布的可见区域。颜色上下限决定数值映射成什么颜色。超过同一上限的多个值可能看起来完全一样，需要分辨大小时，应查看原始空间曲线或数值表。
\subhead{怎样理解两个采样点之间的颜色}
对于物理位置热图，程序使用每个时刻自己的位置向量，将场量线性插值到共同的空间网格。超出该时刻采样支持范围的位置可能留白。时间在横轴时，显示过程还可能沿时间插值。因此，尽管练习原始数据只有 9 个时刻和 21 个中心，热图仍可呈现平滑色带。
对于文件中的质量或编号坐标（\code{CMC}、\code{CMI}、\code{NC}、\code{NI}），同一编号在各时刻有固定坐标。如果文件提供频率群中心及相应边界数据，菜单还可能出现多群 \ui{I-Left} 或 \ui{I-Right}，此时显示的是时间与频率群的关系。
\subhead{查看局部并返回}
在画布中按住鼠标左键拖出矩形，可放大框选区域。滚轮围绕指针位置调整显示范围。先点击画布获得焦点，按 \keys{H} 返回上一级缩放，按 \keys{A} 恢复完整范围。如果正在提取线段或添加标注，先停止相应工具，再拖动缩放。
\checknote{回到原始采样}{平滑色带中间的值可能来自插值。核对细节时，在空间图中选择实际记录的时刻，或在时间图中选定一个具体单元。}
